\documentclass{article}
\usepackage[T1]{fontenc}
\usepackage{amsmath,amsthm,xcolor}
\newenvironment{REV}{\color{red}}{\ignorespacesafterend}
\newenvironment{OLD}{\color{gray}}{\ignorespacesafterend}
\newcommand{\rev}[1]{\textcolor{blue}{#1}}
\newcommand{\phrase}{text generated by a macro}
\newtheorem{theorem}{Theorem}
\begin{document}
\section{Plain text}
MKpa1 A plain paragraph whose source is wrapped by hand, so that a typeset row
usually takes its words from two source lines. MKpa2 sits on the second source line
and MKpa3 on the third, followed by enough words to fill a few rows of the page and
make line breaking matter for the positions of the words. MKpa4 closes it.

Hyphenation of extraordinarily incomprehensibilities and ligatures in fine efficient
offices MKhy1 are typeset here, with an inline formula $a + b
= c$ split over two source lines MKim1 and a citation~\cite{ref} MKci1, then a
footnote mark\footnote{MKfn1 The footnote text sits on this source line.} MKfn0 ends it.
% a comment line MKcm1 that is never typeset
% another comment line MKcm2

\section{Revision proof}
\begin{theorem}\label{thm:main}
MKth1 Every bounded sequence has a convergent subsequence.
\end{theorem}
\begin{REV}
\begin{proof}
MKpr1 We begin the proof with a sentence on its own line.
% a comment inside the paragraph MKcm3
% one more comment line inside the paragraph MKcm4
MKpr2 The paragraph resumes after two comment lines, refers to Theorem~\ref{thm:main}
and to \eqref{eq:one}, and \rev{MKrv1 a revised phrase} continues with \phrase{} MKma1.
\begin{equation}\label{eq:one}
  x = y + z \quad\text{MKeq1}
\end{equation}
MKpr3 Then the desired result follows directly by letting $R \to \infty$ in the display.
\begin{align*}
  u &= v + w \quad\text{MKal1} \\
  v &= w + u \quad\text{MKal2}
\end{align*}
\begin{OLD}
MKol1 An earlier version of this sentence stays visible in gray.
\end{OLD}
MKpr4 This completes the proof.
% \begin{OLD}
% MKcm5 A long commented old block before the end of the proof, as in a revised paper,
% whose lines are long and numerous so that a share of the bytes would land inside it:
% aaaaaaaaaaaaaaaaaaaaaaaaaaaaaaaaaaaaaaaaaaaaaaaaaaaaaaaaaaaaaaaaaaaaaaaaaaaaaaaaaaaaaa
% bbbbbbbbbbbbbbbbbbbbbbbbbbbbbbbbbbbbbbbbbbbbbbbbbbbbbbbbbbbbbbbbbbbbbbbbbbbbbbbbbbbbbb
% cccccccccccccccccccccccccccccccccccccccccccccccccccccccccccccccccccccccccccccccccccccc
% dddddddddddddddddddddddddddddddddddddddddddddddddddddddddddddddddddddddddddddddddddddd
% \end{OLD}
\end{proof}
\end{REV}

\section{Long proof}
\begin{proof}
MKlp01a The long proof continues with paragraph number 1, written over several
source lines so that its rows draw words from more than one line. MKlp01b is on
the second line, and the third line carries MKlp01c together with a short formula
$f(x_1) \le C_1$ and a few more words to finish the paragraph.

MKlp02a The long proof continues with paragraph number 2, written over several
source lines so that its rows draw words from more than one line. MKlp02b is on
the second line, and the third line carries MKlp02c together with a short formula
$f(x_2) \le C_2$ and a few more words to finish the paragraph.

MKlp03a The long proof continues with paragraph number 3, written over several
source lines so that its rows draw words from more than one line. MKlp03b is on
the second line, and the third line carries MKlp03c together with a short formula
$f(x_3) \le C_3$ and a few more words to finish the paragraph.

MKlp04a The long proof continues with paragraph number 4, written over several
source lines so that its rows draw words from more than one line. MKlp04b is on
the second line, and the third line carries MKlp04c together with a short formula
$f(x_4) \le C_4$ and a few more words to finish the paragraph.

MKlp05a The long proof continues with paragraph number 5, written over several
source lines so that its rows draw words from more than one line. MKlp05b is on
the second line, and the third line carries MKlp05c together with a short formula
$f(x_5) \le C_5$ and a few more words to finish the paragraph.

MKlp06a The long proof continues with paragraph number 6, written over several
source lines so that its rows draw words from more than one line. MKlp06b is on
the second line, and the third line carries MKlp06c together with a short formula
$f(x_6) \le C_6$ and a few more words to finish the paragraph.

MKlp07a The long proof continues with paragraph number 7, written over several
source lines so that its rows draw words from more than one line. MKlp07b is on
the second line, and the third line carries MKlp07c together with a short formula
$f(x_7) \le C_7$ and a few more words to finish the paragraph.
% a commented block in the middle of the long proof MKcm6
% second commented line MKcm7
% third commented line MKcm8

MKlp08a The long proof continues with paragraph number 8, written over several
source lines so that its rows draw words from more than one line. MKlp08b is on
the second line, and the third line carries MKlp08c together with a short formula
$f(x_8) \le C_8$ and a few more words to finish the paragraph.

MKlp09a The long proof continues with paragraph number 9, written over several
source lines so that its rows draw words from more than one line. MKlp09b is on
the second line, and the third line carries MKlp09c together with a short formula
$f(x_9) \le C_9$ and a few more words to finish the paragraph.

MKlp10a The long proof continues with paragraph number 10, written over several
source lines so that its rows draw words from more than one line. MKlp10b is on
the second line, and the third line carries MKlp10c together with a short formula
$f(x_10) \le C_10$ and a few more words to finish the paragraph.

MKlp11a The long proof continues with paragraph number 11, written over several
source lines so that its rows draw words from more than one line. MKlp11b is on
the second line, and the third line carries MKlp11c together with a short formula
$f(x_11) \le C_11$ and a few more words to finish the paragraph.

MKlp12a The long proof continues with paragraph number 12, written over several
source lines so that its rows draw words from more than one line. MKlp12b is on
the second line, and the third line carries MKlp12c together with a short formula
$f(x_12) \le C_12$ and a few more words to finish the paragraph.

MKlp13a The long proof continues with paragraph number 13, written over several
source lines so that its rows draw words from more than one line. MKlp13b is on
the second line, and the third line carries MKlp13c together with a short formula
$f(x_13) \le C_13$ and a few more words to finish the paragraph.

MKlp14a The long proof continues with paragraph number 14, written over several
source lines so that its rows draw words from more than one line. MKlp14b is on
the second line, and the third line carries MKlp14c together with a short formula
$f(x_14) \le C_14$ and a few more words to finish the paragraph.

\end{proof}

\section{Included file}
MKch1 Text from an included file, written on the second line of chapter.tex.
\begin{REV}
MKch2 A revised paragraph in the included file
continues with MKch3 on its next source line.
\end{REV}


\begin{thebibliography}{9}
\bibitem{ref} MKbb1 A synthetic reference.
\end{thebibliography}
\end{document}
